% =============================================================================
% CAPÍTULO 7 — INTEGRACIONES Y CONSUMO DE SERVICIOS EXTERNOS
% Archivo maestro del capítulo. Importa las secciones desde sections/cap07/.
% =============================================================================
\chapter{Integraciones y Consumo de Servicios Externos}
\label{ch:integraciones}
\resumenCapitulo{Este capítulo documenta las tres integraciones externas de la plataforma:
Supabase (autenticación, funciones RPC y canales de tiempo real), Google Gemini (motor de
IA generativa) y Google Drive con su infraestructura Picker, detallando el contrato de
cada frontera, sus \textit{payloads}, sus credenciales y sus planes de contingencia.}

EthosHub no es un sistema cerrado: delega capacidades especializadas en \textbf{tres
proveedores externos} cuya integración define buena parte del valor de la plataforma. Este
capítulo documenta —con el detalle que la norma ISO/IEC/IEEE~15289:2024 exige para el ítem
de información de \textit{interfaces externas}— cómo se consume cada servicio, qué datos
cruzan cada frontera y cómo se gobiernan los fallos.

Los capítulos~\ref{ch:backend} (pág.~\pageref{ch:backend}) y~\ref{ch:frontend}
(pág.~\pageref{ch:frontend}) ya describieron \textit{dónde} viven estas integraciones en el
código; este capítulo se centra en el \textbf{contrato de cada frontera}: el protocolo, los
\textit{payloads}, las claves y los planes de contingencia. Las tres integraciones son:

\begin{itemize}[noitemsep]
    \item \textbf{Supabase} (sección~\ref{sec:int_supabase}, pág.~\pageref{sec:int_supabase}):
    autenticación, base de datos, funciones RPC y canales de tiempo real.
    \item \textbf{Google Gemini} (sección~\ref{sec:int_gemini}, pág.~\pageref{sec:int_gemini}):
    motor de IA generativa para CV e \textit{insights} de reclutamiento.
    \item \textbf{Google Drive + Picker} (sección~\ref{sec:int_drive},
    pág.~\pageref{sec:int_drive}): importación de archivos desde el almacenamiento del usuario.
\end{itemize}

La figura~\ref{fig:int-mapa} ofrece un mapa general de las integraciones, distinguiendo
las que consume el \textbf{cliente} directamente de las que median por el \textbf{backend}.

\vspace{0.3cm}
\begin{figure}[h!]
\centering
\begin{tikzpicture}[node distance=0.9cm and 1.7cm,
    box/.style={rectangle, rounded corners=3pt, draw, font=\sffamily\scriptsize, align=center,
        minimum width=2.6cm, minimum height=0.9cm, inner sep=3pt},
    ext/.style={box, fill=c4Externo!30, draw=c4Externo!70!black},
    rel/.style={-{Stealth[length=2mm]}, semithick}]
    \node[box, fill=c4Persona!25] (cli) {Cliente SPA};
    \node[box, fill=c4Componente!45, below=of cli] (be) {Backend\\Spring};
    \node[ext, right=2.2cm of cli] (sb) {Supabase\\\scriptsize Auth · DB · RPC · RT};
    \node[ext, right=2.2cm of be] (gem) {Google\\Gemini};
    \node[ext, below=of gem] (drive) {Google\\Drive + Picker};

    \draw[rel] (cli) -- node[c4etiqueta, above]{\scriptsize anon key + JWT} (sb);
    \draw[rel] (cli) to[bend right=12] node[c4etiqueta, right]{\scriptsize OAuth picker} (drive);
    \draw[rel] (be) -- node[c4etiqueta, above]{\scriptsize service\_role} (sb);
    \draw[rel] (be) -- node[c4etiqueta, above]{\scriptsize API key} (gem);
\end{tikzpicture}
\caption{Mapa general de integraciones externas: vías directas del cliente y mediadas por el backend.}
\label{fig:int-mapa}
\end{figure}

% --- Secciones del capítulo (orden del índice instrucciones.md §7) ---
% =============================================================================
% 7.1 INTEGRACIÓN CON SUPABASE
% =============================================================================
\section{Integración con Supabase}
\label{sec:int_supabase}

Supabase es el proveedor central de EthosHub: aporta la autenticación, la base de datos
PostgreSQL (capítulo~\ref{ch:basedatos}, pág.~\pageref{ch:basedatos}), las funciones RPC y
los canales de tiempo real. La integración es \textbf{dual}: el cliente consume Supabase
\textit{directamente} con la \texttt{anon key} —protegido por las políticas RLS del motor
(sección~\ref{sec:rls}, pág.~\pageref{sec:rls})—, mientras que el backend lo consume con la
\texttt{service\_role key} para operaciones privilegiadas.

\subsection{Cliente del navegador y autenticación}
\label{subsec:supabase_cliente}

El frontend instancia un \textbf{cliente de \texttt{supabase-js}} en \pathbreak{lib/supabase.ts}
con una configuración deliberadamente \textbf{minimalista}: \comando{autoRefreshToken: false},
\comando{persistSession: false} y \comando{detectSessionInUrl: false}. La sesión \textbf{no}
la gestiona Supabase, sino el \comando{authStore} del propio frontend
(capítulo~\ref{ch:frontend}, sección~\ref{subsec:persistencia},
pág.~\pageref{subsec:persistencia}); Supabase se usa solo como \textbf{transporte} de RPC y
tiempo real.

\clearpage
\begin{lstlisting}[language=JavaScript, caption={Cliente de Supabase con inyección manual del JWT (\texttt{lib/supabase.ts}).}, label={lst:supabase-client}]
export const supabase = createClient(supabaseUrl, supabaseAnonKey, {
  auth: { autoRefreshToken: false, persistSession: false,
          detectSessionInUrl: false },
  global: {
    fetch: (url, options = {}) => {            // inyecta el JWT en cada fetch
      const token = sessionStorage.getItem(ACCESS_TOKEN_KEY);
      const headers = new Headers(options.headers);
      if (token) headers.set('Authorization', `Bearer ${token}`);
      return fetch(url, { ...options, headers });
    },
  },
});
\end{lstlisting}

\begin{NotaTecnica}
La bandera \comando{isSupabaseConfigured} comprueba que \texttt{VITE\_SUPABASE\_URL} y
\texttt{VITE\_SUPABASE\_ANON\_KEY} estén presentes antes de crear el cliente; si faltan, el
cliente es \texttt{null} y las funciones de tiempo real degradan con elegancia en lugar de
fallar. Tras autenticarse, \comando{setSupabaseAuth(token)} propaga el JWT al \textit{socket}
de Realtime con \texttt{supabase.realtime.setAuth(...)}, de modo que las suscripciones
respeten igualmente las políticas RLS.
\end{NotaTecnica}

\subsection{Consumo de funciones RPC desde el frontend}
\label{subsec:supabase_rpc}

Para operaciones de negocio que conviene resolver \textbf{atómicamente en el motor}, el
cliente invoca \textbf{funciones RPC} (PL/pgSQL \texttt{SECURITY DEFINER}, documentadas en
el capítulo~\ref{ch:basedatos}, sección~\ref{sec:programable},
pág.~\pageref{sec:programable}) mediante \comando{supabase.rpc(nombre, parámetros)}. La
tabla~\ref{tab:rpcs-fe} cataloga las RPC reales consumidas desde el cliente.

\vspace{0.2cm}
\noindent
\renewcommand{\arraystretch}{1.3}
\fontsize{8.5}{10.2}\selectfont\sffamily
\begin{tabularx}{\textwidth}{|L{4.4cm}|Y|}
\hline
\rowcolor{headerTabla}
\sffamily\textbf{Función RPC} & \sffamily\textbf{Propósito} \\
\hline
\texttt{mark\_messages\_read} & Marca como leídos los mensajes de un chat. \\ \hline
\rowcolor{grisTecnico}
\texttt{get\_candidate\_notes} & Recupera las notas privadas del reclutador sobre un candidato. \\ \hline
\texttt{add\_candidate\_note} · \texttt{update\_candidate\_note} · \texttt{delete\_candidate\_note} & CRUD de notas privadas del reclutador. \\ \hline
\rowcolor{grisTecnico}
\texttt{get\_company\_pipeline} & Devuelve el \textit{pipeline} Kanban de candidatos de la empresa. \\ \hline
\texttt{update\_like\_stage} & Mueve a un candidato de fase en el \textit{pipeline}. \\ \hline
\rowcolor{grisTecnico}
\texttt{remove\_like} & Elimina un \textit{like}/candidato del \textit{pipeline}. \\ \hline
\texttt{upsert\_chat} & Crea o recupera el chat entre reclutador y candidato. \\ \hline
\end{tabularx}
\label{tab:rpcs-fe}

\begin{lstlisting}[language=JavaScript, caption={Invocación de una RPC con parámetros tipados (real).}, label={lst:rpc-call}]
const { data, error } = await supabase.rpc('get_company_pipeline',
  { p_company_id: profile.id });
\end{lstlisting}

\begin{AvisoNota}
El prefijo \texttt{p\_} en los parámetros (\texttt{p\_company\_id}, \texttt{p\_profile\_id})
es la convención de nomenclatura de los argumentos de las funciones PL/pgSQL del esquema.
Como estas RPC son \texttt{SECURITY DEFINER}, ejecutan con privilegios elevados pero validan
internamente la pertenencia, combinando atomicidad con seguridad.
\end{AvisoNota}

\subsection{Canales de comunicación en tiempo real}
\label{subsec:supabase_realtime}

El tiempo real se implementa con el protocolo de \textbf{eventos compartidos} de Supabase,
suscribiéndose a \comando{postgres\_changes}: cambios en tablas del esquema \texttt{core}
que Supabase entrega por \textit{WebSocket} a los clientes suscritos, \textbf{sin
\textit{polling}}. La figura~\ref{fig:rt-canales} ilustra los canales activos.

\vspace{0.3cm}
\begin{figure}[h!]
\centering
\begin{tikzpicture}[node distance=0.7cm and 1.5cm,
    box/.style={rectangle, rounded corners=3pt, draw, font=\sffamily\scriptsize, align=center,
        minimum width=2.7cm, minimum height=0.85cm, inner sep=3pt},
    rel/.style={-{Stealth[length=2mm]}, semithick}]
    \node[box, fill=c4Datos!22, draw=c4Datos!70!black] (db) {core.chat\_messages\\core.profile\_likes\\\scriptsize \dots};
    \node[box, fill=colorAcento!28, draw=colorAcento!80!black, right=of db] (rt) {Realtime\\\scriptsize postgres\_changes};
    \node[box, fill=c4Componente!40, right=of rt, yshift=0.9cm] (chat) {Canal chat\\\scriptsize presencia + mensajes};
    \node[box, fill=c4Componente!40, right=of rt, yshift=-0.9cm] (dash) {Canal dashboard\\\scriptsize métricas en vivo};

    \draw[rel] (db) -- node[c4etiqueta, above]{\scriptsize WAL} (rt);
    \draw[rel] (rt) -- (chat);
    \draw[rel] (rt) -- (dash);
\end{tikzpicture}
\caption{Canales de tiempo real de Supabase sobre cambios en el esquema \texttt{core}.}
\label{fig:rt-canales}
\end{figure}

La tabla~\ref{tab:canales-rt} enumera los canales reales y la tabla que observan.

\vspace{0.2cm}
\noindent
\renewcommand{\arraystretch}{1.3}
\fontsize{8.5}{10.2}\selectfont\sffamily
\begin{tabularx}{\textwidth}{|L{4.4cm}|Y|}
\hline
\rowcolor{headerTabla}
\sffamily\textbf{Canal} & \sffamily\textbf{Tabla observada y efecto} \\
\hline
\texttt{presence:\{chatId\}} & Presencia de los participantes de un chat. \\ \hline
\rowcolor{grisTecnico}
\texttt{global-chat-rec:\{id\}} & \texttt{chat\_messages}: nuevos mensajes para el reclutador. \\ \hline
\texttt{recruiter-dashboard-live:\{id\}} & \texttt{recruiter\_talent\_views} y \texttt{profile\_likes}: refresca métricas. \\ \hline
\end{tabularx}
\label{tab:canales-rt}

\begin{NotaTecnica}
La suscripción a \comando{postgres\_changes} puede dejar de entregar eventos de forma
silenciosa tras una caída de red; por ello las suscripciones inspeccionan el estado del
canal (\texttt{subscribe((status) => ...)}) para detectar la desconexión y re-suscribirse.
Esta resiliencia es la contraparte cliente del modelo de entrega \textit{push} validado por
el backend (capítulo~\ref{ch:backend}, sección~\ref{sec:realtime},
pág.~\pageref{sec:realtime}).
\end{NotaTecnica}
    % 7.1 Integración con Supabase
% =============================================================================
% 7.2 INTEGRACIÓN CON GOOGLE GEMINI AI
% =============================================================================
\section{Integración con Google Gemini AI Service}
\label{sec:int_gemini}

La inteligencia generativa de EthosHub se apoya en \textbf{Google Gemini 2.5 Flash},
consumido \textbf{exclusivamente desde el backend} (nunca desde el cliente, para no exponer
la clave de API). Esta sección documenta el flujo de procesamiento, el mapeo de los
\textit{payloads} hacia el cliente HTTP y el plan de contingencia ante fallos.

\subsection{Flujo de procesamiento asíncrono de IA}
\label{subsec:gemini_flujo}

Gemini atiende dos casos de uso, descritos en el capítulo~\ref{ch:backend}
(sección~\ref{sec:secuencias}, pág.~\pageref{sec:secuencias}): la \textbf{asistencia en
currículums} (\comando{GeminiAiService}, con historial multi-turno) y las \textbf{analíticas
de reclutamiento} (\comando{RecruiterAiService}, con cuatro modos —véase la
tabla~\ref{tab:ia-modos}, pág.~\pageref{tab:ia-modos}). Ambos comparten la misma frontera
HTTP saliente y la misma disciplina de errores. La figura~\ref{fig:gemini-flujo} resume el
procesamiento.

\vspace{0.3cm}
\begin{figure}[h!]
\centering
\begin{tikzpicture}[node distance=0.6cm and 1.2cm,
    box/.style={rectangle, rounded corners=3pt, draw, font=\sffamily\scriptsize, align=center,
        minimum width=2.5cm, minimum height=0.85cm, inner sep=3pt},
    ext/.style={box, fill=c4Externo!30, draw=c4Externo!70!black},
    rel/.style={-{Stealth[length=2mm]}, semithick}]
    \node[box, fill=c4Componente!45] (svc) {Servicio IA\\\scriptsize (CV / reclutador)};
    \node[box, fill=colorAcento!25, right=of svc] (sys) {System\\instruction};
    \node[box, fill=c4Componente!40, below=of sys] (body) {Request body\\\scriptsize JSON + historial};
    \node[ext, right=2.0cm of sys] (api) {Gemini API\\\scriptsize generateContent};
    \node[box, fill=c4Componente!40, below=of api] (parse) {Parse JSON\\\scriptsize extrae texto};

    \draw[rel] (svc) -- node[c4etiqueta, above]{\scriptsize ancla salida} (sys);
    \draw[rel] (sys) -- (body);
    \draw[rel] (body) -- (api);
    \draw[rel] (api) -- node[c4etiqueta, right]{\scriptsize 200 / error} (parse);
    \draw[rel] (parse) to[bend right=30] node[c4etiqueta, below]{\scriptsize sugerencia} (svc);
\end{tikzpicture}
\caption{Flujo de procesamiento de una petición a Gemini desde el backend.}
\label{fig:gemini-flujo}
\end{figure}

\subsection{Mapeo técnico de \textit{payloads} hacia el cliente HTTP}
\label{subsec:gemini_payload}

La comunicación con Gemini usa el \textit{endpoint} \comando{generateContent} del modelo
\texttt{gemini-2.5-flash}. El \textit{payload} se construye con un \textit{system
instruction} (que ancla el comportamiento del modelo) y el contenido del usuario; la
respuesta se navega como árbol JSON para extraer el texto generado. La
tabla~\ref{tab:gemini-payload} mapea los campos del \textit{request}.

\vspace{0.2cm}
\noindent
\renewcommand{\arraystretch}{1.3}
\fontsize{8.5}{10.2}\selectfont\sffamily
\begin{tabularx}{\textwidth}{|L{3.8cm}|Y|}
\hline
\rowcolor{headerTabla}
\sffamily\textbf{Campo del \textit{payload}} & \sffamily\textbf{Contenido} \\
\hline
\texttt{system\_instruction} & Instrucción que fija el rol y el formato de salida (LaTeX seguro o JSON de perfiles). \\ \hline
\rowcolor{grisTecnico}
\texttt{contents} & Historial multi-turno y la petición actual del usuario. \\ \hline
\texttt{generationConfig} & \texttt{maxOutputTokens=8192}, \texttt{temperature=0.7}, \texttt{thinkingBudget=0}. \\ \hline
\end{tabularx}
\label{tab:gemini-payload}

\begin{NotaTecnica}
El transporte se realiza con \comando{HttpURLConnection} (\texttt{connectTimeout}~15\,s,
\texttt{readTimeout}~120\,s). Aunque el índice del proyecto contempla el uso de un cliente
HTTP de mayor nivel (HttpClient5), la implementación actual emplea el cliente nativo de la
JVM por simplicidad y ausencia de dependencias adicionales; el contrato externo (URL,
\textit{payload}, parseo de respuesta) es idéntico e independiente del cliente concreto, de
modo que un cambio de transporte no alteraría la frontera de integración.
\end{NotaTecnica}

\subsection{Plan de contingencia y manejo de errores}
\label{subsec:gemini_contingencia}

La integración con IA es la más frágil del sistema (depende de cuota, disponibilidad del
proveedor y validez de la clave). Por ello el manejo de errores es \textbf{exhaustivo} y
\textbf{nunca propaga un 500 genérico}. La tabla~\ref{tab:gemini-contingencia} documenta
los escenarios de fallo y su tratamiento.

\vspace{0.2cm}
\noindent
\renewcommand{\arraystretch}{1.3}
\fontsize{8.5}{10.2}\selectfont\sffamily
\begin{tabularx}{\textwidth}{|L{3.8cm}|C{1.8cm}|Y|}
\hline
\rowcolor{headerTabla}
\sffamily\textbf{Escenario} & \sffamily\textbf{Respuesta} & \sffamily\textbf{Tratamiento} \\
\hline
Clave de API ausente & 500 controlado & \texttt{IllegalStateException}: ``Asistente IA no configurado''. \\ \hline
\rowcolor{grisTecnico}
Cuota agotada & 429 & Mensaje de límite de uso al usuario. \\ \hline
Proveedor caído (5xx) & 503 & ``El asistente IA no está disponible''. \\ \hline
\rowcolor{grisTecnico}
Otros errores \textit{upstream} & 502 & Código y mensaje del fallo \textit{upstream}. \\ \hline
Respuesta sin texto & 502 & Se registra el cuerpo y se devuelve error accionable. \\ \hline
\end{tabularx}
\label{tab:gemini-contingencia}

\begin{AlertaSeguridad}
La ausencia de \texttt{GEMINI\_API\_KEY} se trata como un \textbf{error de configuración
explícito}, no como un fallo silencioso: el sistema avisa con un mensaje claro
(``añade \texttt{gemini.api-key} en \texttt{application.yml}'') en lugar de enviar
peticiones inválidas. Esta detección temprana —junto al mapeo de estados de la
tabla~\ref{tab:gemini-contingencia}— materializa el principio de degradación elegante: una
funcionalidad de IA caída no derriba el resto de la plataforma.
\end{AlertaSeguridad}
      % 7.2 Integración con Google Gemini
% =============================================================================
% 7.3 INTEGRACIÓN CON GOOGLE DRIVE E INFRAESTRUCTURA PICKER
% =============================================================================
\section{Integración con Google Drive e Infraestructura Picker}
\label{sec:int_drive}

La tercera integración permite al usuario \textbf{importar archivos desde su Google Drive}
sin salir de la aplicación, mediante la \textbf{infraestructura Picker} de Google. A
diferencia de Gemini, este flujo lo gobierna el \textbf{cliente}: es el navegador quien
abre el selector de Google y obtiene el archivo, que luego se canaliza al backend para su
almacenamiento.

\subsection{Arquitectura del flujo de importación desde el frontend}
\label{subsec:drive_flujo}

El flujo combina dos APIs de Google —\textbf{Google Picker API} (el selector visual) y
\textbf{Google Drive API} (el acceso al archivo)— con la subida posterior al \textit{bucket}
de Supabase Storage a través del backend. La figura~\ref{fig:drive-flujo} lo modela.

\vspace{0.3cm}
\begin{figure}[h!]
\centering
\begin{tikzpicture}[node distance=0.65cm and 1.25cm,
    box/.style={rectangle, rounded corners=3pt, draw, font=\sffamily\scriptsize, align=center,
        minimum width=2.5cm, minimum height=0.85cm, inner sep=3pt},
    ext/.style={box, fill=c4Externo!30, draw=c4Externo!70!black},
    rel/.style={-{Stealth[length=2mm]}, semithick}]
    \node[box, fill=c4Persona!25] (user) {Usuario};
    \node[box, fill=c4Componente!45, right=of user] (btn) {Botón\\``Añadir de Drive''};
    \node[ext, right=of btn] (picker) {Google\\Picker API};
    \node[ext, below=of picker] (drive) {Google\\Drive API};
    \node[box, fill=c4Componente!45, below=of btn] (svc) {fileService\\\scriptsize POST /uploads};
    \node[box, fill=c4Datos!22, draw=c4Datos!70!black, left=of svc] (store) {Supabase\\Storage};

    \draw[rel] (user) -- (btn);
    \draw[rel] (btn) -- node[c4etiqueta, above]{\scriptsize abre} (picker);
    \draw[rel] (picker) -- node[c4etiqueta, right]{\scriptsize selección} (drive);
    \draw[rel] (drive) to[bend right=12] node[c4etiqueta, left]{\scriptsize archivo} (svc);
    \draw[rel] (svc) -- node[c4etiqueta, above]{\scriptsize multipart} (store);
\end{tikzpicture}
\caption{Flujo de importación de archivos desde Google Drive vía Picker.}
\label{fig:drive-flujo}
\end{figure}

Una vez seleccionado el archivo en el Picker, su contenido se entrega al
\comando{fileService}, que lo sube por \pathbreak{POST /api/uploads}
(capítulo~\ref{ch:backend}, sección~\ref{subsec:seq_upload},
pág.~\pageref{subsec:seq_upload}); el backend, actuando como intermediario de confianza,
lo persiste en Supabase Storage con la \texttt{service\_role key}. Así, el archivo importado
de Drive sigue exactamente la misma ruta de almacenamiento que cualquier otra subida, sin
caminos privilegiados en el cliente.

\subsection{Tokens y variables de entorno seguras}
\label{subsec:drive_tokens}

La integración requiere credenciales de Google Cloud, declaradas como \textbf{variables de
entorno públicas} del cliente (prefijo \texttt{VITE\_}). La tabla~\ref{tab:drive-vars}
las documenta, tomadas de \pathbreak{.env.example}.

\vspace{0.2cm}
\noindent
\renewcommand{\arraystretch}{1.3}
\fontsize{8.5}{10.2}\selectfont\sffamily
\begin{tabularx}{\textwidth}{|L{4.4cm}|Y|}
\hline
\rowcolor{headerTabla}
\sffamily\textbf{Variable} & \sffamily\textbf{Propósito} \\
\hline
\texttt{VITE\_GOOGLE\_CLIENT\_ID} & ID de cliente OAuth 2.0 (tipo ``Aplicación web'') con el dominio en orígenes autorizados. \\ \hline
\rowcolor{grisTecnico}
\texttt{VITE\_GOOGLE\_API\_KEY} & Clave de API para habilitar Google Picker API y Google Drive API. \\ \hline
\end{tabularx}
\label{tab:drive-vars}

\begin{AlertaSeguridad}
Ambas credenciales son \textbf{públicas por diseño} (viajan en el \textit{bundle} del
cliente), igual que la \texttt{anon key} de Supabase
(sección~\ref{sec:int_supabase}, pág.~\pageref{sec:int_supabase}). Su seguridad no reposa
en el secreto, sino en la \textbf{restricción por origen}: el \texttt{Client ID} de OAuth
solo funciona desde los dominios autorizados en Google Cloud Console, lo que impide su
reutilización desde un origen no aprobado. Las credenciales \textbf{privadas} (como la
\texttt{service\_role} de Supabase o la \texttt{GEMINI\_API\_KEY}) jamás llevan el prefijo
\texttt{VITE\_} y permanecen solo en el backend.
\end{AlertaSeguridad}

\begin{NotaTecnica}
La frontera de Drive es el único punto donde el cliente consume directamente un tercero
distinto de Supabase. El prefijo \texttt{VITE\_} es, en sí mismo, un \textbf{control de
seguridad}: Vite solo expone al \textit{bundle} las variables con ese prefijo, de modo que
un secreto sin él (definido en el entorno del backend) es \textbf{inalcanzable} desde el
navegador por construcción. La gestión integral de secretos se detalla en el
capítulo~\ref{ch:vyv} (sección~\ref{sec:secretos}, pág.~\pageref{sec:secretos}).
\end{NotaTecnica}
       % 7.3 Integración con Google Drive + Picker
